% GENERATED FILE. Edit canonical lessons, then run npm run generate:books.
% canonical-source-hash: fnv1a64:5bfc1a9010445f01
% canonical-lessons: ZH-C461-peifu, ZH-C461-aixin, ZH-W461-ai-love, ZH-C461-aiqing, ZH-C461-xinai, ZH-C461-panwang

\chapter{Love and Kindness}
\label{ch:zh-love-and-kindness}
\hlchaptermodality{461}

\begin{chapteropening}
\textbf{By the end of this chapter:} \emph{I can say whom I admire, talk about kindness, love stories and the things I love, and what I look forward to.}

Complete the last lesson of chapter 461: I can say whom I admire, talk about kindness, love stories and the things I love, and what I look forward to.
\end{chapteropening}

\section[pèifú]{pèifú — to think highly of; to hold in high regard}

\label{lesson:ZH-C461-peifu}



\begin{quote}
Before the new one: say the Mandarin for to sympathize with, then the Mandarin for to envy.
\end{quote}

\subsection*{You'll want to know: pèifú}
\emph{pèifú} — \textquotedblleft{}to think highly of; to hold in high regard\textquotedblright{}.

\begin{quote}
\emph{wǒ zuì pèifú wǒ de lǎoshī} — The person I admire most is my teacher.
\end{quote}

Its characters come later; for now it is a word for the ear.

\begin{cousinweb}
\emph{pèifú} is built from \emph{pèi} 'to wear at the belt' and \emph{fú} 'to serve', both inherited from Middle Chinese, \emph{bwojH} and \emph{bjuwk} (Guangyun, in Baxter's transcription), which glosses the first as jade worn at the belt and the second 'to serve' and 'clothes'. The sense 'to admire' is not given there.
\end{cousinweb}

\begin{grammarlens}[title={What you've built}]
One more word for saying how it was.
\end{grammarlens}

\subsection*{Your turn}
\begin{itemize}
\raggedright
  \item \emph{Say it:} \emph{pèifú}
  \item \emph{Say it:} \emph{pèifú}, once more
  \item \emph{Say it:} \emph{wǒ zuì pèifú wǒ de lǎoshī} — The person I admire most is my teacher.
  \item \emph{Recall:} say the Mandarin for to sympathize with, then the Mandarin for to envy, then say \emph{pèifú} again
\end{itemize}

\begin{tcolorbox}[breakable,colback=teal!4,colframe=teal!35!black,title={Before you move on}]
What does pèifú mean? (\textquotedblleft{}To think highly of; to hold in high regard\textquotedblright{}.) Say it once more.
\end{tcolorbox}
\section[\texorpdfstring{\zh{爱心} (àixīn)}{爱心 (àixīn)}]{\zh{爱心} (àixīn) — compassion; love and care for others}

\label{lesson:ZH-C461-aixin}



\begin{quote}
Before the new one: say the Mandarin for to envy, then the Mandarin for to hold in high regard.
\end{quote}

\subsection*{You'll want to know: \zh{爱心}}
\textbf{\zh{爱心}} — \emph{àixīn} — \textquotedblleft{}compassion; love and care for others\textquotedblright{}.

\begin{quote}
\emph{nàge lùrén hěn yǒu àixīn wǒ hěn gǎnjī tā} — That passer-by was so kind, and I'm very grateful to him.
\end{quote}

The new character, in pieces: \zh{爱} — ten strokes in three pieces: at the top a short stroke falling to the left, two dots and a stroke falling to the left, like a hand reaching down; a cover of two strokes, a short stroke down and a bar that turns down with a hook; and under it a bar, a long stroke falling to the left, and \zh{又} \emph{yòu}. The next lesson puts it in your hand.

\begin{cousinweb}
\textbf{\zh{爱心}} (\emph{àixīn}) is built from \emph{ài} 'to love' and \zh{心} 'heart', both inherited from Middle Chinese, \emph{'ojH} and \emph{sim} (Guangyun, in Baxter's transcription), which glosses the first 'to cherish'.
\end{cousinweb}

\begin{grammarlens}[title={What you've built}]
One more word for saying how it was.
\end{grammarlens}

\subsection*{Your turn}
\begin{itemize}
\raggedright
  \item \emph{Say it:} \emph{àixīn}
  \item \emph{Say it:} \emph{àixīn}, once more
  \item \emph{Say it:} \emph{nàge lùrén hěn yǒu àixīn wǒ hěn gǎnjī tā} — That passer-by was so kind, and I'm very grateful to him.
  \item \emph{Recall:} say the Mandarin for to envy, then the Mandarin for to hold in high regard, then say \emph{àixīn} again
\end{itemize}

\begin{tcolorbox}[breakable,colback=teal!4,colframe=teal!35!black,title={Before you move on}]
What does \zh{爱心} mean? (\textquotedblleft{}Compassion; love and care for others\textquotedblright{}.) Say it once more.
\end{tcolorbox}
\section[\texorpdfstring{\zh{爱} (ài)}{爱 (ài)}]{\zh{爱} — a hand on top, a cover, and \zh{又} under them}

\label{lesson:ZH-W461-ai-love}



\begin{quote}
The word first, then the character inside it.

\begin{itemize}
\raggedright
  \item \emph{Recall:} say \emph{àixīn} — \textbf{R1}, one lesson back, before it decays
\end{itemize}

\begin{quote}
\zh{又} \zh{心}
\end{quote}

\emph{Read it:} each one aloud, with nothing said first — \emph{yòu}, \emph{xīn}
\end{quote}

\subsection*{Script — a hand on top, a cover, and \zh{又} under them}
\hlblockfigure{\detokenize{figures/ZH-W461-ai-love-filmstrip.pdf}}{How \zh{爱} is written, stroke by stroke}

\begin{quote}
\zh{爱}
\end{quote}

\textbf{Three pieces, top to bottom.} Count the pieces, not the strokes.

\textbf{The top, four strokes.} A short stroke from the right falling to the left, a dot, a second dot, then a short stroke falling to the left on the right: a hand reaching down.

\textbf{The cover, two strokes.} A short stroke down on the left, then a bar that runs right, turns down and hooks.

\textbf{The bottom, four strokes.} A bar across, a long stroke falling to the left from above the bar, then \zh{又}: a bar that turns into a stroke falling to the left, and a long stroke falling to the right.

4 + 2 + 4 = ten.

\subsection*{Your turn}
\begin{itemize}
\raggedright
  \item \emph{Write it:} \zh{爱}, all ten, with the model in front of you
  \item \emph{Cover:} write it again from memory, and count the pieces, not the strokes
\end{itemize}

\begin{tcolorbox}[breakable,colback=teal!4,colframe=teal!35!black,title={Before you move on}]
How many pieces does \zh{爱} have? (\textbf{Three: four strokes, two and four}.) Which written character ends it? (\textbf{\zh{又}}.)

Source: \href{https://raw.githubusercontent.com/chanind/hanzi-writer-data/68d10a4b21150cae5e1ebbd223eed289cf32d90c/data/\%E7\%88\%B1.json}{Hanzi Writer Data \zh{爱}, pinned stroke paths and medians}.
\end{tcolorbox}
\section[\texorpdfstring{\zh{爱情} (àiqíng)}{爱情 (àiqíng)}]{\zh{爱情} (àiqíng) — romantic love; love between a couple}

\label{lesson:ZH-C461-aiqing}



\begin{quote}
Before the new one: say the Mandarin for to hold in high regard, then the Mandarin for love and care for others.
\end{quote}

\subsection*{You'll want to know: \zh{爱情}}
\textbf{\zh{爱情}} — \emph{àiqíng} — \textquotedblleft{}romantic love; love between a couple\textquotedblright{}.

\begin{quote}
\emph{zhè shì wǒ kàn guo de zuì hǎo de àiqíng diànyǐng} — This is the best love film I've ever seen.
\end{quote}

\begin{cousinweb}
\textbf{\zh{爱情}} (\emph{àiqíng}) is built from \zh{爱} 'to love' and \zh{情} 'feeling', both inherited from Middle Chinese, \emph{'ojH} and \emph{dzjeng} (Guangyun, in Baxter's transcription), which glosses the first 'to cherish'.
\end{cousinweb}

\begin{grammarlens}[title={What you've built}]
One more word for saying how it was.
\end{grammarlens}

\subsection*{Your turn}
\begin{itemize}
\raggedright
  \item \emph{Say it:} \emph{àiqíng}
  \item \emph{Say it:} \emph{àiqíng}, once more
  \item \emph{Say it:} \emph{zhè shì wǒ kàn guo de zuì hǎo de àiqíng diànyǐng} — This is the best love film I've ever seen.
  \item \emph{Recall:} say the Mandarin for to hold in high regard, then the Mandarin for love and care for others, then say \emph{àiqíng} again
\end{itemize}

\begin{tcolorbox}[breakable,colback=teal!4,colframe=teal!35!black,title={Before you move on}]
What does \zh{爱情} mean? (\textquotedblleft{}Romantic love; love between a couple\textquotedblright{}.) Say it once more.
\end{tcolorbox}
\section[\texorpdfstring{\zh{心爱} (xīn'ài)}{心爱 (xīn'ài)}]{\zh{心爱} (xīn'ài) — beloved; treasured}

\label{lesson:ZH-C461-xinai}



\begin{quote}
Before the new one: say the Mandarin for love and care for others, then the Mandarin for love between a couple.
\end{quote}

\subsection*{You'll want to know: \zh{心爱}}
\textbf{\zh{心爱}} — \emph{xīn'ài} — \textquotedblleft{}beloved; treasured\textquotedblright{}.

\begin{quote}
\textbf{\zh{这是我最心爱的书}} \emph{zhè shì wǒ zuì xīn'ài de shū} — This is my most treasured book.
\end{quote}

\begin{cousinweb}
\textbf{\zh{心爱}} (\emph{xīn'ài}) is built from \zh{心} 'heart' and \zh{爱} 'to love', both inherited from Middle Chinese, \emph{sim} and \emph{'ojH} (Guangyun, in Baxter's transcription), which glosses the second 'to cherish'.
\end{cousinweb}

\begin{grammarlens}[title={What you've built}]
One more word for saying how it was.
\end{grammarlens}

\subsection*{Your turn}
\begin{itemize}
\raggedright
  \item \emph{Say it:} \emph{xīn'ài}
  \item \emph{Say it:} \emph{xīn'ài}, once more
  \item \emph{Say it:} \emph{zhè shì wǒ zuì xīn'ài de shū} — This is my most treasured book.
  \item \emph{Recall:} say the Mandarin for love and care for others, then the Mandarin for love between a couple, then say \emph{xīn'ài} again
\end{itemize}

\begin{tcolorbox}[breakable,colback=teal!4,colframe=teal!35!black,title={Before you move on}]
What does \zh{心爱} mean? (\textquotedblleft{}Beloved; treasured\textquotedblright{}.) Say it once more.
\end{tcolorbox}
\section[pànwàng]{pànwàng — to look forward to; to long for}

\label{lesson:ZH-C461-panwang}



\begin{quote}
Before the new one: say the Mandarin for love between a couple, then the Mandarin for treasured.
\end{quote}

\subsection*{You'll want to know: pànwàng}
\emph{pànwàng} — \textquotedblleft{}to look forward to; to long for\textquotedblright{}.

\begin{quote}
\emph{háizi dōu pànwàng guònián} — All the children are looking forward to the new year.
\end{quote}

Its characters come later; for now it is a word for the ear.

\begin{cousinweb}
\emph{pànwàng} is built from \emph{pàn} and \emph{wàng} 'to look toward', both inherited from Middle Chinese, \emph{pheanH} and \emph{mjangH} (Guangyun, in Baxter's transcription). The dictionary glosses the first as beautiful eyes and the second 'to look toward', after an older book a person far away hoping to come home; the sense 'to look forward to' for the first is not given there.
\end{cousinweb}

\begin{grammarlens}[title={What you've built}]
One more word for saying how it was.
\end{grammarlens}

\subsection*{Your turn}
\begin{itemize}
\raggedright
  \item \emph{Say it:} \emph{pànwàng}
  \item \emph{Say it:} \emph{pànwàng}, once more
  \item \emph{Say it:} \emph{háizi dōu pànwàng guònián} — All the children are looking forward to the new year.
  \item \emph{Recall:} say the Mandarin for love between a couple, then the Mandarin for treasured, then say \emph{pànwàng} again
\end{itemize}

\begin{tcolorbox}[breakable,colback=teal!4,colframe=teal!35!black,title={Before you move on}]
What does pànwàng mean? (\textquotedblleft{}To look forward to; to long for\textquotedblright{}.) Say it once more.
\end{tcolorbox}
